\documentclass[a4paper]{article}     
\usepackage{amsfonts}
\usepackage{amsmath}
\usepackage{amssymb}
\usepackage{graphicx}

\begin{document}

\noindent \large Auteur: Elias Alsa Ati                         \\
\noindent \large Studierichting: Informatica                    \\
\noindent \large Oefeningenreeks 1B nr. 47.1                    \\

\medskip

\normalsize

$$5z + 4 \bar{z} + 5 z^{-1} = 10$$

We vereenvoudigen door $z^{-1}$ weg te werken:  \\

$5z^2 + 4 \bar{z} z + 5 = 10z$                  \\

Berekening van $z^2$:                           \\

$\begin{array}{l}
    z^2 = (a + bi)^2                            \\
    z^2 = (a + bi)(a + bi)                      \\
    z^2 = a^2 - b^2 + 2abi                      \\
\end{array}$                                    \\

Berekening van $z\bar{z}$                       \\

$\begin{array}{l}
    z\bar{z} = (a + bi)(a - bi)                 \\
    z\bar{z} = a^2 - (bi)^2                     \\
    z\bar{z} = a^2 + b^2 = |z|^2                \\
\end{array}$                                    \\

Verder eenvoudiging:                            \\

$\begin{array}{l}
    5\left(a^2 - b^2 + 2abi\right) + 4 \left(a^2 + b^2\right) + 5   = 10 \left(a + bi\right)    \\
    5 a^2 - 5 b^2 + 10 abi + 4 a^2 + 4 b^2 + 5                      = 10 a + 10 bi              \\
    9 a^2 - b^2 + 10 abi + 5                                        = 10 a + 10 bi              \\
\end{array}$                        \\

\textbf{Berekening van stelsel: }   \\

$$\left\{
    \begin{array}{l}
        9a^2 - b^2 + 5 = 10a        \\
        10ab = 10b
    \end{array}
\right .$$

$\begin{array}{l}
    \Rightarrow
    \left\{
        \begin{array}{l}
            9a^2 - b^2 + 5 = 10a    \\
            ab = b
        \end{array}
    \right .
    \Rightarrow
    \left\{
        \begin{array}{l}
            9a^2 - b^2 + 5 = 10a    \\
            ab - b = 0
        \end{array}
    \right .
    \Rightarrow
    \left\{
        \begin{array}{l}
            9a^2 - b^2 + 5 = 10a    \\
            b(a-1) = 0
        \end{array}
    \right .
\end{array}$            \\

\textbf{2 scenario's:}  \\

\textbf{Stel $b = 0$:}  \\

$\begin{array}{l}
    9a^2 - b^2 + 5 = 10a    \\
    9a^2 - 0^2 + 5 = 10a    \\
    9a^2 + 5       = 10a    \\
    9a^2 - 10a + 5 = 0      \\
\end{array}$                \\

We berekenen de discriminant $\alpha = 9$, $\beta = 10$, $\gamma = 5$:      \\

$\begin{array}{l}
    \Delta = \beta^2 - 4 \alpha \gamma  \\
    \Delta = 10^2 - 4 \cdot 9 \cdot 5   \\
    \Delta = 100 - 180                  \\
    \Delta = -80                        \\
\end{array}$  \\

De discriminant is negatief dus zal de uitkomst in $\mathbb{C}$ zitten, maar $a \in \mathbb{R}$.    \\

Dus situatie $b = 0$ klopt niet.                                                                    \\

\textbf{Stel $a = 1$:}  \\

$\begin{array}{l}
    9a^2 - b^2 + 5        = 10a         \\
    9 \cdot 1^2 - b^2 + 5 = 10 \cdot 1  \\
    14 - b^2              = 10          \\
    - b^2                 = - 4         \\
    b                     = \sqrt{4}    \\
    b                     = \pm 2       \\
\end{array}$                            \\

Dus de oplossing voor $z$ zijn:         \\

\begin{eqnarray*}
    1 + 2i, 1 - 2i
\end{eqnarray*}

\end{document} 
